% ============================================================
% CHAPITRE 3 : Systèmes du 1er et 2nd ordre
% ============================================================
\chapterbanner{Chapitre 3 --- Syst\`emes 1\textsuperscript{er} et 2\textsuperscript{nd} ordre}{ch3green}

\begin{tcolorbox}[colback=ch3green!5,colframe=ch3green!80!black,title={Syst\`eme du 1\textsuperscript{er} ordre}]
\tightmath{
\boxed{G(s) = \frac{K}{1 + \tau s}} \quad \text{P\^ole : } s_1 = -1/\tau
}
\textbf{R\'ep. indicielle} (\'echelon $u_0$) :
\tightmath{
\boxed{y(t) = K u_0\!\left(1 - e^{-t/\tau}\right)}
}
\begin{center}
\begin{tikzpicture}[scale=0.44, transform shape, >=Stealth]
  \draw[->] (0,0) -- (5.5,0) node[right]{$t$};
  \draw[->] (0,0) -- (0,2.3) node[above]{$y$};
  \draw[dashed, gray] (0,1.8) -- (5.5,1.8) node[right]{$Ku_0$};
  \draw[thick, ch3green!80!black] (0,0) .. controls (0.8,1.5) and (2,1.75) .. (5.5,1.8);
  \draw[dashed, gray] (1,0) -- (1,1.14); \node[below] at (1,0) {$\tau$};
  \draw[dashed, gray] (2,0) -- (2,1.56); \node[below] at (2,0) {$2\tau$};
  \draw[dashed, gray] (3,0) -- (3,1.71); \node[below] at (3,0) {$3\tau$};
  \draw[red!70, thin, dashed] (0,0) -- (1,1.8);
  \node[above right, red!70] at (0.3,1.8) {tangente};
\end{tikzpicture}
\end{center}
\begin{center}
\renewcommand{\arraystretch}{1}
\scriptsize
\begin{tabular}{@{}cccc@{}}
\toprule
$t$ & $y/(Ku_0)$ & \% & Remarque \\
\midrule
$\tau$ & $0.632$ & 63\% & Tangente coupe $Ku_0$ \\
$2\tau$ & $0.865$ & 86\% & \\
$3\tau$ & $0.950$ & 95\% & $\approx t_s$ \`a 5\% \\
$4\tau$ & $0.982$ & 98\% & $\approx t_s$ \`a 2\% \\
$5\tau$ & $0.993$ & 99\% & \\
\bottomrule
\end{tabular}
\end{center}
\end{tcolorbox}

\begin{tcolorbox}[colback=ch3green!5,colframe=ch3green!80!black,title={1\textsuperscript{er} ordre : rampe \& identification}]
\textbf{R\'eponse \`a rampe} $u(t){=}t$ : $y(t) = K[t - \tau + \tau e^{-t/\tau}]$. Erreur perm. : $e_\infty = K\tau$.

\textbf{Identification :}
\begin{enumerate}
  \item $K = y_\infty / u_0$ (gain statique)
  \item $\tau$ = temps pour 63\% de $y_\infty$, ou tangente \`a l'origine coupe $y_\infty$ en $t{=}\tau$
\end{enumerate}
\end{tcolorbox}

\begin{tcolorbox}[colback=ch3green!5,colframe=ch3green!80!black,title={Syst\`eme du 2\textsuperscript{nd} ordre}]
\tightmath{
\boxed{G(s) = \frac{K\,\omega_n^2}{s^2 + 2\zeta\omega_n s + \omega_n^2}} \quad \boxed{s_{1,2} = -\zeta\omega_n \pm \omega_n\sqrt{\zeta^2 - 1}}
}
\tightmath{
\boxed{s^2 + 2\zeta\omega_n s + \omega_n^2 = (s+\zeta\omega_n)^2 + \omega_n^2(1-\zeta^2)}
}
\tightmath{
\boxed{\omega_n = \sqrt{\delta^2+\omega_0^2}} \quad \boxed{\zeta = \frac{\delta}{\omega_n} = \frac{-\Re(s)}{\omega_n}}
}
\tightmath{
\boxed{G(s)=\frac{b_0}{s^2+as+b}\ \Rightarrow\ \omega_n=\sqrt{b},\ \zeta=\frac{a}{2\sqrt{b}},\ K=\frac{b_0}{b}}
}
\begin{itemize}
  \item $\omega_n$ : pulsation naturelle non amortie (rad/s)
  \item $\zeta$ : taux d'amortissement (sans unit\'e)
  \item $K$ : gain statique ($= G(0)$)
\end{itemize}
\begin{center}
\renewcommand{\arraystretch}{1}
\scriptsize
\begin{tabular}{@{}llll@{}}
\toprule
$\zeta$ & R\'egime & P\^oles & Comportement \\
\midrule
$0$ & Oscillant pur & $\pm j\omega_n$ & Oscil. perp\'etuelle \\
$0{<}\zeta{<}1$ & \textbf{Sous-amorti} & $-\delta \pm j\omega_0$ & Oscil. amorties \\
$1$ & Critique & $-\omega_n$ (dble) & + rapide sans $D_\%$ \\
$>1$ & Sur-amorti & 2 r\'eels n\'eg. & Lent, pas de $D_\%$ \\
\bottomrule
\end{tabular}
\end{center}
\end{tcolorbox}

\begin{tcolorbox}[colback=ch3green!5,colframe=ch3green!80!black,title={Cas sous-amorti ($0 < \zeta < 1$) : formules cl\'es}]
\tightmath{
\boxed{\omega_0 = \omega_n\sqrt{1{-}\zeta^2}} \quad \delta = \zeta\omega_n \quad \varphi = \arccos(\zeta)
}
R\'eponse indicielle :
\tightmath{
y(t) = K\!\left[1 - \frac{e^{-\delta t}}{\sqrt{1{-}\zeta^2}}\sin(\omega_0 t + \varphi)\right]
}
\tightmath{
\boxed{D_\% = e^{-\pi\zeta/\sqrt{1-\zeta^2}} \times 100\%} \quad \boxed{\zeta = \frac{-\ln(D_\%/100)}{\sqrt{\pi^2 + \ln^2(D_\%/100)}}}
}
\tightmath{
\boxed{t_p = \frac{\pi}{\omega_0}} \quad \boxed{t_{s2} \approx \frac{4}{\zeta\omega_n}} \quad \boxed{t_{s5} \approx \frac{3}{\zeta\omega_n}} \quad t_r \approx \frac{1.8}{\omega_n}
}
\piege{$\zeta \approx 0.7 \Rightarrow D_\% \approx 5\%$ : bon compromis rapidit\'e/stabilit\'e. Table $\zeta\!\to\!D_\%$ d\'etaill\'ee dans l'abaque ci-dessous.}
\end{tcolorbox}

\begin{tcolorbox}[colback=ch3green!5,colframe=ch3green!80!black,title={Position des p\^oles dans le plan $s$}]
\begin{center}
\begin{tikzpicture}[scale=0.82, transform shape, >=Stealth]
  \draw[->] (-2.8,0) -- (0.8,0) node[right]{Re};
  \draw[->] (0,-2.2) -- (0,2.2) node[above]{Im};
  \filldraw[ch3green!80!black] (-1.2,1.5) circle (2pt) node[right]{$s_1$};
  \filldraw[ch3green!80!black] (-1.2,-1.5) circle (2pt) node[right]{$s_1^*$};
  \draw[dashed, gray, thin] (0,0) -- (-1.2,1.5);
  \draw[<->, thin, red!70] (-1.2,0) -- (-1.2,1.5) node[midway, right]{$\omega_0$};
  \draw[<->, thin, red!70] (0,0) -- (-1.2,0) node[midway, below]{$\delta$};
  \node at (-0.3,1.0) {$\omega_n$};
  \draw[thin] (-0.35,0) arc (180:128:0.35); \node at (-0.55,0.28) {$\theta$};
\end{tikzpicture}
\end{center}
\end{tcolorbox}

\begin{tcolorbox}[colback=ch3green!5,colframe=ch3green!80!black,title={Cas critique, sur-amorti \& identification}]
\textbf{Critique ($\zeta{=}1$)} : $y(t) = K[1-(1{+}\omega_n t)e^{-\omega_n t}]$. $t_s \approx 4.75/\omega_n$.

\textbf{Sur-amorti ($\zeta{>}1$)} : 2 p\^oles r\'eels, $\tau_{1,2} = -1/s_{1,2}$.
\tightmath{
G(s) = \frac{K}{(1{+}\tau_1 s)(1{+}\tau_2 s)}
}
\textbf{Identification 2\textsuperscript{nd} ordre :}
\begin{enumerate}
  \item $K = y_\infty/u_0$
  \item $D_\% \Rightarrow \zeta$ (formule inverse)
  \item $t_p \Rightarrow \omega_0 = \pi/t_p \Rightarrow \omega_n = \omega_0/\sqrt{1{-}\zeta^2}$
\end{enumerate}
\piege{Z\'ero demi-plan droit ($\tau_z{<}0$) $\Rightarrow$ r\'ep. inverse initiale = non-minimum de phase.}
\end{tcolorbox}

\begin{tcolorbox}[colback=ch3green!5,colframe=ch3green!75!black,title={Approximation par p\^ole dominant}]
\textbf{P\^ole dominant} = p\^ole r\'eel n\'egatif \textbf{le plus proche de 0} :
\tightmath{
\tau_i = \frac{-1}{s_i} = \frac{1}{|s_i|} \qquad s_i < 0
}
\vars{Plus $|s_i|$ est petit $\Rightarrow$ plus $\tau_i$ est grand $\Rightarrow$ plus lent $\Rightarrow$ domine.}
\vars{Si un p\^ole est en $s=0$, alors c'est lui le dominant : il est plus lent que tout p\^ole de partie r\'eelle strictement n\'egative.}

\textbf{Ex.} : $s_1{=}{-}100$,\ $s_2{=}{-}20$,\ $s_3{=}{-}2$ :
\begin{center}
\renewcommand{\arraystretch}{1.0}\scriptsize
\begin{tabular}{@{}ccc@{}}
\toprule
$s_i$ & $\tau_i$ & \\
\midrule
$-100$ & $0.01$\,s & n\'egligeable \\
$-20$  & $0.05$\,s & n\'egligeable \\
$\boldsymbol{-2}$ & $\boldsymbol{0.5}$\,\textbf{s} & \textbf{dominant} \\
\bottomrule
\end{tabular}
\end{center}
\tightmath{
G(s) \approx \frac{K_{\text{dom}}}{1 + 0.5\,s}
}

\textbf{Annulation p\^ole-z\'ero} : si $s_z \approx s_p$ (tous deux $< 0$) :
\tightmath{
\frac{1{+}\tau_z s}{1{+}\tau_p s} \approx 1
}
\piege{Z\'ero en demi-plan \textbf{droit} ($s_z > 0$) : pas de simplification.}
\end{tcolorbox}

